% =====================================================================
% preamble.tex — パッケージ読込・色・カスタム環境・マクロの定義
% 本文（chapters/*.tex）から使う見た目の部品はすべてここに集約する。
% エンジン: LuaLaTeX + jlreq
% =====================================================================

% --- 基本パッケージ ---------------------------------------------------
\usepackage{graphicx}          % 画像の取り込み
\usepackage{xcolor}            % 色指定
\usepackage{booktabs}          % 罫線のきれいな表
\renewcommand{\arraystretch}{1.25} % 表の行を少しゆったりさせる
\usepackage{tabularx}          % 幅を自動調整する表
\usepackage{enumitem}          % 箇条書きの調整
\usepackage{fancyvrb}          % コード枠に使う Verbatim
\usepackage[most]{tcolorbox}   % 装飾枠（コード・注意書き）
\usepackage{adjustbox}         % 図が本文幅を超えたときの自動縮小
\usepackage{tikz}              % 概念図・フロー図
\usetikzlibrary{arrows.meta,positioning,shapes.geometric,calc,fit,backgrounds}
\usepackage{float}             % 図を本文の位置に固定する [H]
\usepackage{caption}           % 図番号つきキャプションの体裁
\captionsetup{font=small, labelfont=bf, labelsep=quad}

% 画像の探索パス（既存の Markdown 版と同じ images/ を共有する）
\graphicspath{{../images/}{./images/}}

% --- 配色 -------------------------------------------------------------
% Markdown 版の Mermaid 図の色をそのまま踏襲する（薄い色は使わない方針）。
\definecolor{gitRed}{HTML}{C0392B}     % 作業ツリー・注意
\definecolor{gitOrange}{HTML}{D68910}  % ステージング
\definecolor{gitGreen}{HTML}{2E7D32}   % リポジトリ・完了
\definecolor{gitBlue}{HTML}{4A90D9}    % コミット・情報
\definecolor{gitPurple}{HTML}{6C3483}  % 自分側
\definecolor{gitTeal}{HTML}{117A65}    % 両方採用
\definecolor{gitAmber}{HTML}{8A6D0B}   % 判断ノード
\definecolor{codeBg}{HTML}{F4F5F7}     % コード枠の背景
\definecolor{codeFrame}{HTML}{D0D3D9}  % コード枠の枠線

% --- コード枠 ---------------------------------------------------------
% 使い方: \begin{codebox}[bash] ... Verbatim ... \end{codebox}
% 省略時のラベルは bash。右上に言語ラベルを小さく表示する。
\newtcolorbox{codebox}[1][bash]{%
  enhanced, breakable,
  colback=codeBg, colframe=codeFrame, boxrule=0.6pt, arc=2pt,
  left=9pt, right=9pt, top=13pt, bottom=8pt,
  before skip=8pt, after skip=8pt,
  overlay={%
    \node[anchor=north east, font=\sffamily\scriptsize\bfseries,
          text=gray!60] at ([xshift=-6pt, yshift=-3pt]frame.north east) {#1};%
  },
}
% コード本体を書くための Verbatim 環境（小さめ・行間やや詰め）
\DefineVerbatimEnvironment{code}{Verbatim}{fontsize=\small, xleftmargin=0pt}

% 本文中のインラインコード: \code{git status}
\newcommand{\cmd}[1]{\texttt{#1}}

% --- 注意書き（admonition）枠 ----------------------------------------
% ヒント（緑）・注意（赤）・補足（グレー）の3種。emoji は使わずラベルで表す。
% タイトルは本文（明るい背景）に直接乗るため、白文字だと見えない。
% タイトル文字色は各テーマ色にして、はっきり読めるようにする。
\newtcolorbox{hintbox}{%
  enhanced, breakable,
  colback=gitGreen!5, colframe=gitGreen, boxrule=0.8pt, arc=2pt,
  left=8pt, right=8pt, top=6pt, bottom=6pt,
  coltitle=gitGreen, fonttitle=\sffamily\bfseries\small,
  title=ヒント, attach title to upper=\quad,
}
\newtcolorbox{warnbox}{%
  enhanced, breakable,
  colback=gitRed!5, colframe=gitRed, boxrule=0.8pt, arc=2pt,
  left=8pt, right=8pt, top=6pt, bottom=6pt,
  coltitle=gitRed, fonttitle=\sffamily\bfseries\small,
  title=注意, attach title to upper=\quad,
}
% 執筆メモ（雛形用の枠）。原稿を清書したら、この枠ごと削除する。
\newtcolorbox{draftbox}{%
  enhanced, breakable,
  colback=gitAmber!6, colframe=white, boxrule=0pt, arc=2pt,
  borderline={0.9pt}{0pt}{gitAmber, dashed},
  left=8pt, right=8pt, top=6pt, bottom=6pt,
  coltitle=gitAmber, fonttitle=\sffamily\bfseries\small,
  title=執筆メモ, attach title to upper=\quad,
}
\newtcolorbox{notebox}{%
  enhanced, breakable,
  colback=gray!8, colframe=gray!60!black, boxrule=0.8pt, arc=2pt,
  left=8pt, right=8pt, top=6pt, bottom=6pt,
  coltitle=gray!30!black, fonttitle=\sffamily\bfseries\small,
  title=補足, attach title to upper=\quad,
}

% --- スクリーンショット取り込み --------------------------------------
% \screenshot[幅の割合]{ファイル}{キャプション}
% フロートにせず本文の流れに沿って中央寄せで置く（説明の直後に見せたいため）。
\newcommand{\screenshot}[3][0.7]{%
  \begin{center}
    \includegraphics[width=#1\linewidth]{#2}\\[3pt]
    {\small\itshape #3}%
  \end{center}%
}

% \screenshottrim[幅の割合]{ファイル}{キャプション}{trim指定}
% 余白の多いスクリーンショットを切り詰めて大きく見せる。
% trim は「左 下 右 上」の順で、画像1ピクセル = 1bp として指定する
% （例: 1920x1080 の下半分を捨てる → {0 540bp 0 0}）。
\newcommand{\screenshottrim}[4][0.95]{%
  \begin{center}
    \includegraphics[width=#1\linewidth, trim=#4, clip]{#2}\\[3pt]
    {\small\itshape #3}%
  \end{center}%
}

% --- TikZ 共通スタイル ------------------------------------------------
% Mermaid のフロー図を再現するためのノード／矢印スタイル。
\tikzset{
  flownode/.style={%
    rounded corners=4pt, draw=none, text=white,
    font=\sffamily\small, align=center,
    minimum height=1.15cm, minimum width=2.2cm, inner sep=6pt},
  decisionnode/.style={%
    diamond, aspect=1.9, draw=none, text=white,
    font=\sffamily\small, align=center, inner sep=3pt},
  flowarrow/.style={-{Latex[length=3.6mm, width=3mm]}, line width=1.8pt, draw=black},
  edgelabel/.style={font=\sffamily\normalsize\bfseries, fill=white, inner sep=3pt},
}
% 図全体を中央寄せで置くためのヘルパ環境（図番号なし）。
% 本文幅を超えた図は自動で縮小する（幅が足りている図は拡大しない）。
\newenvironment{diagram}%
  {\par\medskip\begin{center}\begin{adjustbox}{max width=\linewidth}}%
  {\end{adjustbox}\end{center}\par\medskip}

% --- 図番号つきの図（本文が「図1.2のように」と参照する概念図・フロー図用）---
% 手順に紐づく操作スクショには使わない（そちらは \screenshot を使う）。
\makeatletter
% TikZ など本体を囲む環境: \begin{numfig}{キャプション}{ラベル} ... \end{numfig}
\newenvironment{numfig}[2]{%
  \def\nf@cap{#1}\def\nf@lab{#2}%
  \begin{figure}[H]\centering\begin{adjustbox}{max width=\linewidth}%
}{%
  \end{adjustbox}\caption{\nf@cap}\label{\nf@lab}\end{figure}%
}
\makeatother
% 画像1枚を図番号つきで置く: \figimage[幅]{ファイル}{キャプション}{ラベル}
\newcommand{\figimage}[4][0.95]{%
  \begin{figure}[H]\centering
    \includegraphics[width=#1\linewidth]{#2}
    \caption{#3}\label{#4}
  \end{figure}%
}
